\slajdid{02-s047}
\begin{frame}[label=02-s047]
Isto tako algebarskim manipulacijama, na račun povećanog kašnjenja mogli smo realizaciju u obliku zbira proizvoda da uradimo i na sledeći način
\[
F=\bar D\bar B+D\bar C\bar A=\bar D\bar B+(D\bar C)\bar A
\]
pa bi nam trebao jedan čip sa invertorima jedan čip sa dvoulaznim I kolima (treba nam ukupno tri i ima ih u jednom čipu) i jedan sa dvoulaznim ILI kolima što je ukupno 3.
\end{frame}
